% --- 1. INFORMAZIONI GENERALI ---
\section{Informazioni Generali}
\begin{itemize}
    \item \textbf{Luogo/Piattaforma:} Server Discord
    \item \textbf{Data:} 2026-03-13
    \item \textbf{Orario di inizio:} 14:30
    \item \textbf{Orario di fine:} 16:20
    \item \textbf{Responsabile:} Francesco Marcon
    \item \textbf{Scriba:} Sofia De Blasi
\end{itemize}

\vspace{0.5cm}
\textbf{Partecipanti presenti:}
\begin{table}[ht]
    \centering
    \renewcommand{\arraystretch}{1.2}
    \begin{tabularx}{\textwidth}{X l}
        \toprule
        \textbf{Nome e Cognome} & \textbf{Matricola} \\
        \midrule
        Sofia De Blasi & 2111014 \\
        Roberto Mariano Doroftei & 2111031 \\
        Filippo Idiometri & 2035617 \\
        Francesco Marcon & 2110990 \\
        Armando Moda Scarati & 2082864 \\
        Anton Ricardo Rupi & 2054304 \\
        \bottomrule
    \end{tabularx}
\end{table}

\vspace{0.5cm}
\textbf{Partecipanti assenti:}
\begin{table}[ht]
    \centering
    \renewcommand{\arraystretch}{1.2}
    \begin{tabularx}{\textwidth}{X l}
        \toprule
        \textbf{Nome e Cognome} & \textbf{Matricola} \\
        \midrule
        Nessuno & - \\
        \bottomrule
    \end{tabularx}
\end{table}

% --- 2. ORDINE DEL GIORNO ---
\section{Ordine del Giorno}
Gli argomenti previsti per la riunione odierna sono:
\begin{enumerate}
    \item Scelta della tecnologia per il tracciamento dei task.
    \item Definizione delle misurazioni: tempi, costi e strumenti di raccolta.
    \item Convenzioni di nomenclatura per file, branch e glossario.
    \item Pianificazione dell’invio delle email alle aziende.
    \item Assegnazione delle attività operative per candidatura e documentazione.
\end{enumerate}

% --- 3. RESOCONTO E DISCUSSIONE ---
\section{Resoconto della Riunione}

\subsection{Configurazione GitHub Project e misurazioni}
Il gruppo ha discusso la necessità di adottare un sistema strutturato per il tracciamento delle attività e la raccolta delle misurazioni.

\textbf{Motivazione:} garantire un processo di gestione dei task coerente, misurabile e integrato con gli strumenti utilizzati dal gruppo.

Azioni operative:
\begin{itemize}
    \item Adozione di GitHub Project come strumento principale per il tracciamento dei task.
    \item Ogni task avrà un identificativo GitHub univoco, utilizzato come riferimento per le misurazioni.
    \item Le misurazioni (ruolo, tempo previsto, tempo effettivo) saranno raccolte tramite un Google Form collegato a un foglio Excel condiviso.
    \item Lo storico delle issue rimarrà su GitHub; le misurazioni saranno gestite tramite Excel.
    \item Tutti i membri leggeranno la documentazione ufficiale di GitHub Project.
\end{itemize}

\subsection{Convenzioni di nomenclatura}
\textbf{Motivazione:} garantire uniformità e coerenza nella struttura della documentazione e nell’organizzazione del repository.

Convenzioni stabilite:
\begin{itemize}
    \item Utilizzo del pedice \textbf{G} per i termini presenti nel glossario.
    \item Struttura della documentazione nella repository:
    \begin{itemize}
        \item \texttt{/docs/CC}, \texttt{/docs/RTB}, \texttt{/docs/PB}.
        \item \texttt{/verb\_interni/verb\_int\_AAAA\_MM\_GG}.
        \item \texttt{/verb\_esterni/verb\_est\_AAAA\_MM\_GG}.
    \end{itemize}
    \item Ogni membro utilizzerà un branch personale.
\end{itemize}

\subsection{Regole formali per decisioni e task}
\textbf{Motivazione:} garantire tracciabilità, coerenza e uniformità nella gestione delle decisioni e delle attività operative all’interno dei verbali.

Convenzioni stabilite:
\begin{itemize}
    \item Le decisioni dei verbali interni saranno identificate tramite il codice \textbf{VI-y.x}, dove:
    \begin{itemize}
        \item \textbf{VI} indica \textit{Verbale Interno};
        \item \textbf{y} rappresenta il numero della riunione;
        \item \textbf{x} è il numero progressivo della decisione.
    \end{itemize}

    \item Le attività operative (To-Do) saranno identificate tramite il codice \textbf{T-y.z}, dove:
    \begin{itemize}
        \item \textbf{T} indica \textit{Task};
        \item \textbf{y} rappresenta il numero della riunione;
        \item \textbf{z} è il numero progressivo del task.
    \end{itemize}

    \item L’ordinamento delle attività nel To-Do seguirà il criterio:
    \begin{enumerate}
        \item scadenza crescente;
        \item decisione di riferimento;
        \item codice incrementale.
    \end{enumerate}
\end{itemize}

\subsection{Comunicazioni con le aziende}
\textbf{Motivazione:} ottenere informazioni aggiuntive sui capitolati e fissare eventuali incontri.

Azioni operative:
\begin{itemize}
    \item Zucchetti: completamento considerazioni tecniche (assegnato a Francesco), invio email (assegnato a Filippo).
    \item Vargroup: invio email (assegnato ad Armando).
\end{itemize}

\subsection{Assegnazione attività operative}
\textbf{Motivazione:} distribuire in modo chiaro e tracciabile le attività necessarie per la candidatura.

Azioni operative:
\begin{itemize}
    \item Stesura del verbale: Sofia.
    \item Creazione del template per i task: Francesco.
    \item Configurazione GitHub Project: Roberto.
    \item Studio delle Norme di Progetto (ISO): Filippo, Roberto, Sofia.
    \item Analisi del preventivo costi: Armando, Ricardo, Francesco.
    \item Studio della documentazione necessaria per la candidatura: tutti.
\end{itemize}

% --- 4. DECISIONI PRESE ---
\section{Decisioni Prese}

\renewcommand{\arraystretch}{1.3}
\vspace{-0.3cm}
\noindent
\begin{longtable}{c p{12cm}}
    \toprule
    \textbf{Codice} & \textbf{Decisione} \\
    \midrule
    \endfirsthead

    \toprule
    \textbf{Codice} & \textbf{Decisione} \\
    \midrule
    \endhead

    \midrule
    \multicolumn{2}{r}{\textit{Continua nella pagina successiva}} \\
    \midrule
    \endfoot

    \bottomrule
    \endlastfoot

    \textbf{VI-2.1} & Adozione di GitHub Project come strumento principale per il tracciamento dei task. \\
    \textbf{VI-2.2} & Raccolta delle misurazioni tramite Excel con eventuale Google Form collegato. \\
    \textbf{VI-2.3} & Definizione delle convenzioni di nomenclatura: pedice G, cartelle interne/esterne, branch per persona. \\
    \textbf{VI-2.4} & Adozione delle regole formali per l’identificazione di decisioni (VI-y.x) e task (T-y.z). \\
    \textbf{VI-2.5} & Invio email alle aziende Zucchetti ed Eggon per richiesta informazioni. \\
    \textbf{VI-2.6} & Assegnazione dei compiti operativi per candidatura e documentazione. \\
\end{longtable}

% --- 5. AZIONI (TODO) ---
\section{Azioni da Svolgere (To-Do)}

\renewcommand{\arraystretch}{1.3}

\begin{longtable}{c c p{4.5cm} l l}

        \toprule
    \textbf{Codice} & \textbf{Rif.} & \textbf{Descrizione Task} & \textbf{Assegnatario} & \textbf{Scadenza} \\
    \midrule
    \endfirsthead

    \toprule
    \textbf{Codice} & \textbf{Rif.} & \textbf{Descrizione Task} & \textbf{Assegnatario} & \textbf{Scadenza} \\
    \midrule
    \endhead

    \midrule
    \multicolumn{5}{r}{\textit{Continua nella pagina successiva}} \\
    \midrule
    \endfoot

    \bottomrule
    \endlastfoot

    % --- 16 MARZO ---
    \textbf{T-2.1} & VI-2.5 & Invio email a Zucchetti e Vargroup. & Filippo & 2026-03-16 \\
    
    % --- 17 MARZO ---
    \textbf{T-2.2} & VI-2.1 & Lettura documentazione GitHub Project. & Tutti & 2026-03-17 \\
    \textbf{T-2.3} & VI-2.1 & Configurazione GitHub Project. & Roberto & 2026-03-17 \\
    \textbf{T-2.4} & VI-2.2 & Creazione foglio Excel per misurazioni e valutazione Google Form. & Tutti & 2026-03-17 \\
    \textbf{T-2.5} & VI-2.6 & Stesura del verbale della riunione del 2026-03-13. & Sofia & 2026-03-17 \\
    \textbf{T-2.6} & VI-2.6 & Creazione template form per task. & Francesco & 2026-03-17 \\
    \textbf{T-2.7} & VI-2.6 & Studio Norme di Progetto (ISO 12207 - 1995). & Filippo, Roberto, Sofia & 2026-03-17 \\
    \textbf{T-2.8} & VI-2.6 & Analisi preventivo costi. & Armando, Ricardo, Francesco & 2026-03-17 \\
    \textbf{T-2.9} & VI-2.6 & Studio documentazione per candidatura. & Tutti & 2026-03-17 \\
\end{longtable}

% --- 6. PROSSIMA RIUNIONE ---
\section{Prossima Riunione}
\begin{itemize}
    \item \textbf{Data:} 2026-03-17
    \item \textbf{Ora:} 14:30
    \item \textbf{OdG preliminare:} Revisione GitHub Project, verifica misurazioni, avanzamento candidatura con inizio stesura ``Valutazione capitolati'', ``Norme di Progetto'' e ``Preventivo dei Costi e Assegnazione dei Ruoli''.
\end{itemize}
